\chapter{光和光的传播}
\problem{1-1}{太阳与月球的直径分别是 $1.39\times10^6\,\mathrm{km}$ 和 $3.5\times10^3\,\mathrm{km}$，太阳到地面的距离为 $1.50\times10^8\,\mathrm{km}$，月球到地面的距离为 $3.8\times10^5\,\mathrm{km}$。试计算，地面上能见到日全蚀区域的面积（可把该区域的地面视为平面）。}
\answerarea
\problemfig[0.34]{1-2}{由立方体的玻璃切下一角做成的棱镜，称为隅角棱镜（见题图）。证明从斜面射入的光线经其它三面全反射后，出射线的方向总与入射线相反。设想一下，这样的棱镜可以在什么场合发挥作用。}{1_2}
\answerarea
\clearpage
\problemfig[0.38]{1-3}{光线射入如题图所示的棱镜，经两次折射和反射后射出。
\begin{enumerate}\item 证明偏向角与入射方向无关，恒等于 $2\alpha$。\item 在此情况下，能否产生色散？\end{enumerate}}{1_3}
\answerarea
\problem{1-4}{试证明：当一束光线通过平行平面玻璃板时，出射光线方向不变，但产生侧向平移。当入射角 $\theta$ 很小时，位移为
\[x=\frac{n-1}{n}\theta t,\]
式中 $n$ 为玻璃板的折射率，$t$ 为其厚度。}
\answerarea
\clearpage
\problem{1-5}{证明：光线相继经过几个平行分界面的多层介质时，出射光线的方向只与两边的折射率有关，与中间各层介质无关。}
\answerarea
\problem{1-6}{顶角 $\alpha$ 很小的棱镜称为光楔（wedge）。证明光楔使垂直入射的光线产生偏向角 $\delta=(n-1)\alpha$，其中 $n$ 是光楔的折射率。}
\answerarea
\clearpage
\problem{1-7}{本题图所示是一种求折射线方向的简捷作图法。例如为了求光线通过棱镜的路径（图 b），可如图 a 以 $O$ 为中心作二圆弧，半径正比于折射率 $n$、$n'$（设 $n>n'$）。作 $OR$ 平行于入射线 $DE$，作 $RP$ 平行于棱镜第一界面的法线 $N_1N'_1$，则 $OP$ 的方向即为第一次折射后光线 $EF$ 的方向。再作 $QP$ 平行于第二界面的法线 $N_2N'_2$，则 $OQ$ 的方向即为出射线 $FG$ 的方向，从而 $\angle ROQ=\delta$ 为偏向角。试论证此法的依据。\fig[0.82]{1_7}}
\answerarea
\problem{1-8}{利用上题的图证明最小偏向角的存在及（1.23）式。
\begin{equation*}
n=\frac{\sin\dfrac{\alpha+\delta_{\min}}{2}}{\sin\dfrac{\alpha}{2}}.
\tag{1.23}
\end{equation*}
其中 $\alpha$ 为棱镜顶角，$\delta_{\min}$ 为最小偏向角。}
\answerarea
\clearpage
\problem{1-9}{已知棱镜顶角为 $60^\circ$，测得最小偏向角为 $53^\circ14'$，求棱镜的折射率。}
\answerarea
\problem{1-10}{顶角为 $50^\circ$ 的三棱镜的最小偏向角是 $35^\circ$，如果把它浸入水中，最小偏向角等于多少？（水的折射率为 $1.33$。）}
\answerarea
\clearpage
\problemfig[0.32]{1-11}{如本题图所示，在水中有两条平行光线1和2，光线2射到水和平行平板玻璃的分界面上。
\begin{enumerate}\item 两光线射到空气中是否还平行？\item 如果光线1发生全反射，光线2能否进入空气？\end{enumerate}}{1_11}
\answerarea
\problem{1-12}{计算光在下列介质之间穿行时的全反射临界角：（1）从玻璃到空气，（2）从水到空气，（3）从玻璃到水。}
\answerarea
\clearpage
\problem{1-13}{设光导纤维玻璃芯和外套的折射率分别为 $n_1$ 和 $n_2$（$n_1>n_2$），垂直端面外介质的折射率为 $n_0$（见题图）。试证明，能使光线在纤维内发生全反射的入射光束的最大孔径角 $\theta_1$ 满足下式：
\[n_0\sin\theta_1=\sqrt{n_1^2-n_2^2}.\]
（$\theta_1$ 称为纤维的数值孔径。）\fig[0.7]{1_13}}
\answerarea
\problem{1-14}{光导纤维外套由折射率为 $1.52$ 的冕牌玻璃做成，芯线由折射率为 $1.66$ 的火石玻璃做成，求垂直端面的数值孔径。}
\answerarea
\clearpage
\problem{1-15}{极限法测液体折射率的装置如本题图所示，$ABC$ 是直角棱镜，其折射率 $n_g$ 为已知。将待测液体涂一薄层于其上表面 $AB$，覆盖一块毛玻璃，用扩展光源在掠入射的方向照明。从棱镜的 $AC$ 面出射的光线的折射角将有一下限 $i'$（用望远镜观察，则在视场中出现有明显分界线的全明半暗区）。试证明，待测液体的折射率 $n$ 可按下式算出：
\[n=\sqrt{n_g^2-\sin^2 i'}.\]
用这种方法测液体折射率，测量范围受到什么限制？\fig{1_15}}
\answerarea
\problem{1-16}{在空气中钠黄光的波长为 $589.3\,\mathrm{nm}$，问：（1）其频率有多大？（2）在折射率为 $1.52$ 的玻璃中其波长为多少？}
\answerarea
\clearpage
\problem{1-17}{在熔融石英中波长为 $550.0\,\mathrm{nm}$ 的光频率为多少？已知折射率为 $1.460$。}
\answerarea
\problem{1-18}{填充下表中的空白：
\begin{center}\renewcommand{\arraystretch}{1.5}\begin{tabular}{|c|c|c|c|c|}\hline
谱线&\multicolumn{2}{c|}{F线}&\multicolumn{2}{c|}{D线}\\\hline
介质&真空&水&真空&水\\\hline
折射率&&1.337&&1.333\\\hline
波长&$486.1\,\mathrm{nm}$&&$589.3\,\mathrm{nm}$&\\\hline
频率&&&&\\\hline 光速&&&&\\\hline
\end{tabular}\end{center}}
\answerarea
\clearpage
\problem{1-19}{一只船以速度 $v$ 在水面上行驶，$v$ 比水面波的速率 $u$ 大。用惠更斯作图法证明，在水中出现一圆锥形波前，其半顶角 $\theta$ 由下式给出：
\[\sin\theta=\frac{u}{v}.\]
此波称为船波，超音速飞机在空气中产生的冲击波，也是这样产生的。设想一下，若电子以大于介质中光速的速度在介质中作匀速运动，会产生什么现象？}
\answerarea
\problem{1-20}{证明图1-18中光线 $A_1C_1$、$A_2B_2C_2$、$A_3B_3C_3$、$\cdots$、$A_nB_n$ 的光程相等。
\bookfig[0.82]{1_18}{图1-18\quad 用惠更斯原理解释反射定律和折射定律}}
\answerarea
\clearpage
\problem{1-21}{证明图1-18中光线 $A_1D_1$、$A_2B_2D_2$、$A_3B_3D_3$、$\cdots$、$A_nB_n$ 的光程相等。
\bookfig[0.82]{1_18}{图1-18\quad 用惠更斯原理解释反射定律和折射定律}}
\answerarea
\problemfig[0.37]{1-22}{在离桌面 $1.0\,\mathrm m$ 处有盏 $100\,\mathrm{cd}$ 的电灯 $L$，设 $L$ 可看作是各向同性的，求桌上 $A$、$B$ 两点的照度（见本题图）。}{1_22}
\answerarea
\clearpage
\problem{1-23}{若上题中电灯 $L$ 可垂直上下移动，问怎样的高度使 $B$ 点的照度最大。}
\answerarea
\problem{1-24}{（1）设天空为亮度均匀的朗伯体，其亮度为 $B$，试证明，在露天水平面上的照度为 $E=\pi B$；（2）在上面的计算中，与我们假设天空是怎样形状的发光面有无关系？与被照射面的位置有无关系？}
\answerarea
\clearpage
\problem{1-25}{试证明，一个理想漫射体受到照度为 $E$ 的辐射时，反射光的亮度 $B=E/\pi$。}
\answerarea
\problem{1-26}{阳光垂直照射地面时，照度为 $1.0\times10^5\,\mathrm{lx}$。若认为太阳的亮度与光流方向无关，并忽略大气对光的吸收，且已知地球轨道半径为 $1.5\times10^8\,\mathrm{km}$，太阳的直径为 $1.4\times10^6\,\mathrm{km}$，求太阳的亮度。}
\answerarea
\clearpage
